% Generated draft: requires mathematical and source review.
\documentclass[11pt]{article}
\usepackage[margin=25mm]{geometry}
\usepackage{fontspec}
\setmainfont{Latin Modern Roman}
\usepackage{amsmath,amssymb,mathtools,unicode-math}
\setmathfont{Latin Modern Math}
\usepackage{enumitem,needspace,xurl,hyperref}
\setlength{\parindent}{0pt}
\setlength{\parskip}{5pt}
\setlength{\emergencystretch}{3em}
\newcommand{\N}{\mathbb{N}}
\newcommand{\Z}{\mathbb{Z}}
\newcommand{\Q}{\mathbb{Q}}
\newcommand{\R}{\mathbb{R}}
\newcommand{\C}{\mathbb{C}}
\newcommand{\F}{\mathbb{F}}
\newcommand{\A}{\mathbb{A}}
\newcommand{\PP}{\mathbb{P}}
\newcommand{\E}{\mathbb{E}}
\newcommand{\eps}{\varepsilon}
\newcommand{\dd}{\,\mathrm{d}}
\begin{document}
\section*{UnsolvedMath ID: 30000151}
\textbf{Source problem code:} catalogue entry
\subsection{Definitions and mathematical statement}
A standard periodic elliptic framework takes an operator
\[
 L=\sum_{|\alpha|\leq2m}a_\alpha(x)\partial^\alpha
\]
on \(\R^n\), with coefficients periodic under a full-rank lattice and, for a system, matrix-valued coefficients acting on \(\C^r\). The exact coefficient regularity, ellipticity convention, reality/self-adjointness, and lower-order assumptions are not supplied by the source. Second-order scalar divergence form is one starting subclass, not an assumption on all requested extensions.

One common Liouville property in this setting is finite dimensionality of every polynomial-growth solution space
\[
 \mathcal V_N(L)=\left\{u:Lu=0,\ \sup_x(1+|x|)^{-N}\|u(x)\|<\infty\right\},
 \qquad\dim\mathcal V_N(L)<\infty\quad(N\in\Z_{\geq0}),
\]
with the appropriate classical or weak solution class. The record does not explicitly define ``the Liouville theorem,'' so this convention must be checked against the original context. It is not the same assertion as saying every bounded solution is constant; constants need not even solve an operator with arbitrary lower-order terms.

The corrected source asks the following linked questions:
\begin{enumerate}
\item Which Liouville results extend beyond second-order scalar divergence-form equations to other periodic scalar equations, higher-order operators, and matrix-valued systems?
\item What criteria decide whether a given periodic elliptic equation has the intended property?
\item Can the decision be effective? An algorithmic version must specify how the coefficients and hypotheses are encoded, with a finite or computable input model; a qualitative criterion alone does not supply an algorithm.
\item Can analogous criteria be obtained for elliptic operators on abelian coverings? A standard model is a regular covering of a compact base with deck group \(\Z^d\), an invariant operator, and polynomial growth measured in a deck word metric or a compatible lifted metric. The exact cover and growth conventions remain to be fixed.
\end{enumerate}
The question does not demand a single unsupported theorem covering every operator in these classes, nor a particular formula for the dimensions.

The supplied literature triage dated 21 August 2026 reports necessary-and-sufficient Liouville criteria by Kuchment--Pinchover for broad periodic elliptic classes and systems, but says the extracted request extends beyond those stated results. The actual theorem hypotheses and criterion are not reproduced, so no universal higher-order/matrix criterion or effective procedure is asserted here. The record retains its supplied partial-resolution assessment.
\subsection{Short English statement}
The source seeks Liouville criteria beyond scalar divergence-form equations, including higher-order systems and abelian covers, and asks for effective decisions. The growth-based property, operator hypotheses, and algorithmic input model need to be specified before a single theorem can be stated.
\subsection{Sources}
\begin{itemize}
\item Local UnsolvedMath registry, permanent ID \(30000151\), exact source code \texttt{catalogue entry}, Oberwolfach Reports Open Problems. Statement status \texttt{corrected\_verified}; the supplied verification reports restoration of all extension and effective-decision questions from the official report.
\item ``Spectral Analysis of Partial Differential Equations (2005),'' Oberwolfach Reports, pp. 2839--2912, DOI \texttt{10.4171/owr/2004/53}: \url{https://doi.org/10.4171/owr/2004/53}. Proposed year and DOI refer to 2004, while the citation label says 2005.
\item Supplied ``Kuchment--Pinchover'': \url{https://doi.org/10.1090/S0002-9947-07-04196-7}. The other supplied literature source is the original report; full theorem hypotheses and criterion are not reproduced.
\item Literature checked 21 August 2026; supplied status \texttt{partially\_solved}. Registry created 25 August 2026 and updated 16 September 2026. This dated curation triage requires independent verification before citation.
\item Registry curation attribution: Ulam AI / UnsolvedMath Contributors, CC BY 4.0. Underlying report and literature retain their supplied provenance; this statement-only draft requires human review.
\end{itemize}
\end{document}
